% Generated draft: requires mathematical and source review.
\documentclass[11pt]{article}
\usepackage[margin=25mm]{geometry}
\usepackage{fontspec}
\setmainfont{Latin Modern Roman}
\usepackage{amsmath,amssymb,mathtools,unicode-math}
\setmathfont{Latin Modern Math}
\usepackage{enumitem,needspace,xurl,hyperref}
\setlength{\parindent}{0pt}
\setlength{\parskip}{5pt}
\setlength{\emergencystretch}{3em}
\newcommand{\N}{\mathbb{N}}
\newcommand{\Z}{\mathbb{Z}}
\newcommand{\Q}{\mathbb{Q}}
\newcommand{\R}{\mathbb{R}}
\newcommand{\C}{\mathbb{C}}
\newcommand{\F}{\mathbb{F}}
\newcommand{\A}{\mathbb{A}}
\newcommand{\PP}{\mathbb{P}}
\newcommand{\E}{\mathbb{E}}
\newcommand{\eps}{\varepsilon}
\newcommand{\dd}{\,\mathrm{d}}
\begin{document}
\section*{UnsolvedMath ID: 701118}
\textbf{Source problem code:} catalogue entry
\subsection{Definitions and mathematical statement}
Let \(p\) be a prime. A finite \(p\)-group is a finite group of order \(p^n\) for some integer \(n\geq 0\). For a group \(G\), its Schur multiplier is the abelian group \(M(G)=H_2(G,\Z)\), where the coefficients have trivial \(G\)-action.

Is the following statement true? For every prime \(p\) and every finite abelian \(p\)-group \(A\), there exists a nonabelian finite \(p\)-group \(G\) such that
\[
M(G)\cong A.
\]
Here the prime is the same for \(A\) and \(G\), and nonabelian means that there exist \(x,y\in G\) with \(xy\ne yx\). Under the usual convention above, the trivial group is included among the possible groups \(A\); the supplied statement gives no exception for it.
\subsection{Short English statement}
Must every finite abelian \(p\)-group occur, up to isomorphism, as the Schur multiplier of some nonabelian finite group whose order is a power of the same prime \(p\)?
\subsection{Sources}
\begin{itemize}
\item Local UnsolvedMath registry: permanent ID \(701118\), exact source problem code \texttt{catalogue entry}, title \emph{Kourovka Notebook Problem 15.30}. The supplied registry labels the problem as open; this status has not been independently verified.
\item \emph{The Kourovka Notebook (Unsolved Problems in Group Theory)}, 15th issue (2002), Problem 15.30. Proposed by Ya. G. Berkovich, according to the supplied record.
\item The supplied background also cites the 21st-issue edition (2026), \texttt{arXiv:1401.0300v48}. The registry describes the collection as covering Issues 1--21 (1965--2026), edited by E. I. Khukhro and V. D. Mazurov. These bibliographic details are supplied attribution and have not been independently checked.
\end{itemize}
\end{document}
